%! AP Statistics — Unit 2, Lesson 2.6 — Exit Ticket — Answer Key
\documentclass[10pt]{article}
\usepackage{apstats-article}
\usepackage{apstats-key}

\begin{document}

\pageheader{Unit 2, Lesson 2.6}{Exit Ticket --- Answer Key}
\namedateperiod

\vspace{0.22in}

% ── SCENARIO ─────────────────────────────────────────────────────────────────
\begin{tcolorbox}[colback=sky, colframe=navy, boxrule=0.9pt, arc=2mm,
    left=4mm, right=4mm, top=3mm, bottom=3mm]
\small\textit{A statistics teacher recorded the number of hours each student studied for a
unit exam and the score each student earned (out of 100). Using data from 28 students who
studied between 0 and 8 hours, technology produced the following regression equation:}
\[
  \widehat{\text{score}} = 40 + 5(\text{hours})
\]
\end{tcolorbox}

\vspace{0.2in}

\begin{enumerate}[leftmargin=1.6em, itemsep=0.2in]

  \item Interpret the \textbf{slope} ($5$) in context.\\[10pt]
        \ansline{For each additional hour studied, the predicted exam score increases by
        approximately 5 points.}

  \item Interpret the \textbf{$y$-intercept} ($40$) in context.
        State whether it is meaningful in this context.\\[10pt]
        \ansline{When a student studies 0 hours, the predicted exam score is 40. This
        \textbf{is meaningful in context} because studying 0 hours is within the data range
        (0--8 hr) and is plausible.}

  \item Predict the exam score for a student who studied \textbf{6 hours}.
        Is this interpolation or extrapolation? Show your calculation.\\[10pt]
        \ansline{$\widehat{\text{score}} = 40 + 5(6) = 40 + 30 = 70$ points. This is
        \textbf{interpolation} because 6 hours is within the observed range of 0--8 hours,
        so the prediction is reliable.}

  \item A student claims she studied \textbf{30 hours}. Use the equation to compute a
        predicted score. Then explain why this prediction is not trustworthy.\\[10pt]
        \ansline{$\widehat{\text{score}} = 40 + 5(30) = 40 + 150 = 190$ points.
        This is \textbf{extrapolation}: 30 hours is far outside the data range (0--8 hr).
        The prediction of 190 is impossible (above 100\%) and illustrates that the linear
        model should not be used outside the range of the observed data.}

\end{enumerate}

\vspace{0.16in}

% ── AP-STYLE MC ──────────────────────────────────────────────────────────────
\begin{tcolorbox}[colback=redbg, colframe=redacc, boxrule=0.8pt, arc=2mm,
    left=4mm, right=4mm, top=3mm, bottom=3mm]
\small\textbf{AP-Style Practice}\\[6pt]
A researcher fits a regression line for data on $x$ = age of a tree (years) and $y$ = trunk
diameter (cm). The data were collected from trees aged 5 to 30 years. The slope is $1.8$ and
the $y$-intercept is $3.5$. Which of the following is the best interpretation of the slope?

\vspace{10pt}
\begin{enumerate}[label=\textbf{(\Alph*)}, leftmargin=2em, itemsep=6pt]
  \item The predicted trunk diameter of a newly planted tree (age $= 0$) is 3.5 cm.
  \item For every additional centimeter of trunk diameter, the predicted age increases by 1.8 years.
  \item For each additional year of age, the predicted trunk diameter increases by approximately 1.8 cm.
        \textcolor{keyred}{\textbf{$\leftarrow$ correct}}
  \item A tree that is 1.8 years old will have a trunk diameter of 3.5 cm.
\end{enumerate}
\end{tcolorbox}

\vspace{0.1in}
\noindent\textcolor{slate}{\small My answer: \ans{C} \quad Because: \ans{(C) correctly gives direction, magnitude, units, and context for the slope; (A) describes the intercept; (B) swaps $x$ and $y$; (D) confuses coefficients.}}

\begin{teachernote}
Use exit ticket results diagnostically. Students who miss \#1--2 need practice with the AP
template. Students who miss \#4 (extrapolation) should revisit the notes example where the
prediction for age 20 gave a negative car price. The MC distractor (B) --- swapping
explanatory and response --- is the most common error; address it explicitly.
\end{teachernote}

\end{document}
